\paragraph{Data in brief.}
The unit of analysis is the calendar day. For each of the four cities
the benchmark provides a daily panel: the count of 311 requests created
in city $c$ and harmonized service family $s$ on day $t$, joined with
deterministic calendar information and one station's weather
observations. Section~4 details sources, exclusions, harmonization, and
provenance; everything below is defined on this panel.

\paragraph{Prediction task.}
\emph{Input: the panel through day $t$. Output: next-day forecasts for
every active family. Goal: measure what each public signal adds, under
strict leakage control.}
Let $y_{c,s,t}$ denote the number of 311 requests created in city $c$ and
harmonized service family $s \in \mathcal{S}_c$ on calendar day $t$, where
$\mathcal{S}_c$ is the city's \emph{active} family set (eight families in
New York, San Francisco, and Austin; seven in Chicago, where residential
noise is structurally absent from the intake taxonomy). At the end of day
$t$ (the information cutoff), a forecaster produces a point forecast
$\hat{y}_{c,s,t+1}$ and, for the probabilistic model, quantile forecasts
$\hat{q}^{\,\tau}_{c,s,t+1}$, $\tau \in \{.05,.25,.50,.75,.95\}$.
Endogenous features use observations through day $t$ only; calendar
features of day $t{+}1$ are deterministic; target-day weather uses the
realized observation for $t{+}1$ as a stand-in for a day-ahead weather
forecast (assumption A3, bounded by a lagged-weather variant).

\paragraph{Why a simulated allocation layer.}
Accuracy metrics alone cannot reveal whether differences between
forecasts would ever change a decision. The benchmark therefore couples
every forecaster to one deliberately transparent allocation simulation:
a scientific instrument for measuring the forecast-to-decision mapping
under a known objective, with every parameter exposed as configuration
and every allocation step auditable. The instrument models no city's
actual operations; its value is that any difference in simulated outcome
is attributable, by construction, to the forecasts and the policy rule
alone.

\paragraph{Simulated allocation.}
\emph{Input: the day-$t{+}1$ forecasts and the observable carryover
state. Output: an integer split of a fixed capacity budget across
families. Goal: score each forecast by the unresolved workload its
allocations leave behind.}
Each evening a simulated planner allocates a fixed budget of $B_c$
abstract capacity units across active families,
$x_{c,s,t+1} \in \mathbb{Z}_{\ge 0}$, $\sum_s x_{c,s,t+1} = B_c$; a unit
resolves up to $\kappa$ request-equivalents the next day. Unresolved
workload carries over:
\begin{align}
  w_{c,s,t+1} &= b_{c,s,t} + y_{c,s,t+1}, \nonumber \\
  \mathrm{resolved} &= \min(w_{c,s,t+1},\, \kappa\, x_{c,s,t+1}), \nonumber \\
  u_{c,s,t+1} &= w_{c,s,t+1} - \mathrm{resolved}, \nonumber \\
  b_{c,s,t+1} &= (1-\alpha)\, u_{c,s,t+1},
\end{align}
where $b$ is simulated unresolved request-equivalent carryover (observable
state at decision time), initialized at $b_{c,s,0}=0$, and $\alpha$ an
abandonment rate ($\alpha=0$ primary; $b_0$ and $\alpha$ are varied as
labeled sensitivities). In words: each day's workload $w$ is new requests
plus carried stock; allocated capacity resolves up to $\kappa$
request-equivalents per unit; whatever remains unresolved ($u$) carries
into the next day's state $b$. The decision loss over the test horizon
$\mathcal{T}$ is
\begin{equation}
  L_c \;=\; \sum_{t \in \mathcal{T}} \sum_{s \in \mathcal{S}_c}
  \omega_s\, u_{c,s,t},
  \label{eq:decision-loss}
\end{equation}
with EQUAL request-equivalent weights $\omega_s = 1$ as the primary
objective --- a deliberately simple equal-weight choice, not a normatively
neutral or normative priority policy (one illustrative normative-priority
scenario is a labeled sensitivity).
Because $u_{c,s,t}$ contains the carried stock, a request that remains
unresolved for $k$ days contributes $k$ times to
\eqref{eq:decision-loss}: the loss is a \emph{waiting-time-weighted
cumulative unresolved stock}, not a daily flow. Under the scarce regime
the simulated system is structurally under-capacitated (supercritical),
so this stock grows over the horizon and daily losses --- and the daily
loss \emph{differentials} between policies --- can carry a deterministic
trend; Section~8 reports a trend diagnostic and the inference framing
follows it.
Hypothetical service-pressure regimes set $\kappa B_c$ to 70\%, 90\%, and
110\% of the city's TRAINING-window mean daily demand $\bar d_c$: that is,
$B_c=\max\!\big(\mathrm{round}(f\,\bar d_c/\kappa),\,|\mathcal S_c|\big)$
for $f\in\{0.7,0.9,1.1\}$, floored at the active-family count. The
resulting integer budgets are frozen before any model selection and used
unchanged for validation-stage selection and the single test evaluation;
$\kappa$ sets only the scale of a capacity unit. Rather than sweeping a
global $\kappa$ grid, the sensitivity analysis varies one family's
service-yield multiplier at a time (factors $0.70$ and $1.30$). Nothing in
this construction estimates any city's actual capacity, productivity, or
queues.

\paragraph{Policies.}
\emph{Input: the same observable state for every policy. Output: a full
allocation of the budget. Goal: isolate how the forecast's
representation, not its information set, changes the allocation.}
All policies receive the same observable state and allocate the full
budget: \emph{uniform} (equal units; forecast-free floor);
\emph{proportional} (largest-remainder apportionment of
$\hat{y}_{c,s,t+1} + b_{c,s,t}$); and \emph{greedy expected-value}, which
exactly maximizes $\sum_s \omega_s\, \mathbb{E}[\min(D_s, \kappa x_s)]$
over integer allocations, where $D_s$ is the forecast
demand-plus-carryover distribution --- degenerate for point forecasts,
piecewise-linear from the quantile model. Greedy is optimal for this
objective because $\mathbb{E}[\min(D, c)]$ is concave in $c$; ties in
marginal value are resolved by a fixed secondary key (ascending family
index). Under a degenerate point forecast this objective is
\emph{non-identified} across most of the feasible set, so the realized
allocation --- and hence its loss --- depends on that key; we treat the
key as a first-class object and report its effect in Section~7. A
hindsight myopic reference (greedy on realized demand) appears in
supplementary diagnostics only; being per-day myopic it is not a
horizon-optimal bound and is never used as a denominator or target.

\paragraph{Controlled uncertainty contrast.}
\emph{Input: one fitted quantile model. Output: three policy arms that
differ only in the representation handed to the same policy. Goal:
isolate the value of distributional information.}
The value of distributional information is isolated within ONE fitted
quantile model: a degenerate-median arm, an implied-mean arm, and a
full-distribution arm differ only in how the same predictive distribution
enters the greedy policy. The predictive distribution is reconstructed by
linear interpolation of the inverse CDF through the five forecast
quantiles on a fixed 99-point grid on $[0.01,0.99]$, with flat (clamped)
tails beyond the $.05$ and $.95$ levels; the implied mean is the average
of the 99 grid samples and the median arm is the $.50$ quantile. We write
``quantile-interpolated predictive distribution'' for this object
throughout. The separately trained point models serve as
forecasting benchmarks, not as the uncertainty comparator.

\paragraph{Evaluation questions.}
For each city and regime we compare forecast accuracy (MAE, with paired
28-day moving-block bootstrap intervals; pinball loss and coverage for
quantiles) and simulated decision loss \eqref{eq:decision-loss} (raw
losses, absolute paired daily differences with block-bootstrap intervals,
and percentage reduction relative to the uniform floor), asking when
predictive and decision orderings agree, whether distributional
information changes outcomes, and whether validation-time selection by
decision loss differs from selection by accuracy.
